\phantomsection
\section*{Chapitre 4 : Mise en \oe{}uvre empirique et r\'esultats}
\addcontentsline{toc}{section}{Chapitre 4 : Mise en \oe{}uvre empirique et r\'esultats}

\bigskip

Le chapitre~3 a pos\'e le cadre m\'ethodologique~: cinq mod\`eles parcimonieux
en concurrence sur chaque ligne budg\'etaire, deux d\'etecteurs d'anomalies
(score~Z univari\'e et Isolation Forest multivari\'e), articulation finale
dans une matrice de priorit\'e P0--P4. Ce chapitre passe \`a la mise en
\oe{}uvre. Il d\'ecrit les donn\'ees effectivement disponibles au Minist\`ere
de la Justice, le protocole exp\'erimental retenu pour discriminer les
mod\`eles, puis pr\'esente les r\'esultats num\'eriques obtenus sur les
107~lignes budg\'etaires communes aux exercices 2021--2026.

Le fil conducteur est double~: chaque chiffre report\'e ici doit \^etre
\emph{reproductible} \`a partir des fichiers \texttt{SituationChap-AAAA.xls}
livr\'es chaque ann\'ee par le minist\`ere, et chaque indicateur agr\'eg\'e
doit pouvoir \^etre \emph{remont\'e \`a la ligne} pour justifier une d\'ecision
op\'erationnelle. La section~4.1 d\'ecrit la cha\^ine d'extraction~; la~4.2
pr\'esente le jeu de donn\'ees consolid\'e et ses caract\'eristiques
statistiques~; la~4.3 d\'etaille le protocole de backtest~; la~4.4 expose
les r\'esultats de s\'election des mod\`eles ligne par ligne~; la~4.5
pr\'esente la pr\'evision triennale 2026--2028~; la~4.6 documente la
d\'etection d'anomalies et la matrice P0--P4~; la~4.7 conclut par une
synth\`ese critique avant le passage au d\'eploiement op\'erationnel
(chapitre~5).

\bigskip

\noindent\textbf{D\'emarche m\'ethodologique en cinq \'etapes.} Avant de
d\'etailler les r\'esultats, il est utile de rappeler l'enchainement logique
qui a structur\'e le travail empirique~:

\medskip

\noindent\textbf{1. D\'efinition du probl\`eme.} La probl\'ematique \'etudi\'ee
consiste \`a pr\'evoir l'\'evolution des lignes budg\'etaires \`a partir d'un
nombre limit\'e d'observations annuelles. Dans le contexte du Minist\`ere de la
Justice, l'objectif est de produire des pr\'evisions \`a trois ans, conform\'ement
\`a la logique de la programmation budg\'etaire triennale instaur\'ee par la
LOF~130-13. En compl\'ement de cette pr\'evision, l'\'etude vise \'egalement \`a
d\'etecter les anomalies susceptibles de signaler des comportements
budg\'etaires atypiques n\'ecessitant une attention particuli\`ere du
gestionnaire.

\medskip

\noindent\textbf{2. Collecte des informations.} Les donn\'ees utilis\'ees
proviennent de deux sources compl\'ementaires. D'une part, des
\textbf{donn\'ees budg\'etaires} fournies par les responsables de la Direction du
Budget ont \'et\'e mobilis\'ees~: bien qu'elles ne correspondent pas exactement
aux donn\'ees r\'eelles d'ex\'ecution -- pour des raisons de confidentialit\'e
institutionnelle -- elles reproduisent fid\`element les caract\'eristiques
structurelles du syst\`eme budg\'etaire \'etudi\'e (nomenclature, ordres de
grandeur, dynamique pluriannuelle). L'objectif principal n'\'etant pas de
reproduire une situation sp\'ecifique mais de concevoir un syst\`eme d'aide \`a
la d\'ecision \emph{g\'en\'eralisable}, ces donn\'ees sont parfaitement
adapt\'ees \`a la construction et \`a l'\'evaluation du pipeline propos\'e.
D'autre part, l'\textbf{expertise m\'etier} des gestionnaires budg\'etaires a
constitu\'e une source d'information essentielle~: les responsables
interrog\'es ont notamment soulign\'e que les donn\'ees trop anciennes peuvent
perdre de leur pertinence en raison de changements structurels affectant
l'environnement budg\'etaire (r\'eformes administratives, inflation,
\'ev\'enements exceptionnels comme la pand\'emie de COVID-19). Cette
observation a conduit \`a privil\'egier les donn\'ees les plus r\'ecentes lors
de la phase de mod\'elisation.

\medskip

\noindent\textbf{3. Analyse exploratoire des donn\'ees.} Avant toute
mod\'elisation, une analyse exploratoire a \'et\'e r\'ealis\'ee afin de mieux
comprendre les caract\'eristiques du ph\'enom\`ene \'etudi\'e. Cette \'etape,
synth\'etis\'ee dans la page de vue d'ensemble du tableau de bord Power~BI
(chapitre~5, section~5.4.1), vise \`a fournir une lecture globale du portefeuille
avant d'aborder la pr\'evision ou la d\'etection d'anomalies. L'analyse a
notamment port\'e sur l'identification de tendances g\'en\'erales, de
comportements r\'ecurrents, de ruptures \'eventuelles ainsi que de valeurs
atypiques. Elle permet \'egalement de confronter les observations statistiques
\`a l'expertise m\'etier afin d'interpr\'eter correctement les anomalies
d\'etect\'ees et d'\'eviter les conclusions erron\'ees~-- principe rappel\'e en
section~3.5.3 (\og un mod\`ele n'est pas un substitut au jugement professionnel \fg{}).

\medskip

\noindent\textbf{4. Choix et ajustement des mod\`eles.} Le choix des mod\`eles
d\'epend de plusieurs facteurs~: la quantit\'e de donn\'ees historiques
disponibles, les caract\'eristiques des s\'eries budg\'etaires \'etudi\'ees et
les objectifs op\'erationnels du syst\`eme de pilotage. Dans le cadre de cette
\'etude, cinq approches de pr\'evision parcimonieuses (justifi\'ees au
chapitre~3) ont \'et\'e mises en concurrence afin d'identifier les mod\`eles
les plus adapt\'es \`a des s\'eries annuelles courtes. Parall\`element, deux
m\'ethodes de d\'etection d'anomalies ont \'et\'e s\'electionn\'ees -- score~Z
univari\'e et Isolation Forest multivari\'e -- afin de rep\'erer les lignes
budg\'etaires au comportement inhabituel. Chaque mod\`ele repose sur un
ensemble d'hypoth\`eses explicites et n\'ecessite une phase d'ajustement \`a
partir des donn\'ees historiques disponibles.

\medskip

\noindent\textbf{5. Utilisation et \'evaluation des mod\`eles.} Une fois les
mod\`eles calibr\'es, ils sont utilis\'es pour g\'en\'erer des pr\'evisions
budg\'etaires et identifier les anomalies potentielles. La qualit\'e des
pr\'evisions est \'evalu\'ee par le MAPE de validation, m\'etrique adapt\'ee \`a
la mesure de l'\'ecart relatif entre valeurs observ\'ees et valeurs pr\'edites
(section~3.5.1). La d\'etection d'anomalies, quant \`a elle, est \'evalu\'ee
sur la capacit\'e des m\'ethodes retenues \`a mettre en \'evidence des
comportements atypiques \emph{coh\'erents avec l'expertise m\'etier}. Les
r\'esultats obtenus sont enfin int\'egr\'es dans un syst\`eme d'aide \`a la
d\'ecision (matrice P0--P4) permettant aux gestionnaires d'identifier
rapidement les lignes \`a surveiller en priorit\'e et de soutenir le pilotage
budg\'etaire op\'erationnel.

\medskip

Les sections suivantes pr\'esentent ces cinq \'etapes sur les donn\'ees du
Minist\`ere de la Justice.

%================================================================
\phantomsection
\subsection*{4.1 \quad Construction du jeu de donn\'ees}
\addcontentsline{toc}{subsection}{4.1 Construction du jeu de donn\'ees}
%================================================================

\subsubsection*{4.1.1 \quad Sources et p\'erim\`etre}
\addcontentsline{toc}{subsubsection}{4.1.1 \quad Sources et p\'erim\`etre}

La source primaire est constitu\'ee des fichiers \texttt{SituationChap-AAAA.xls}
produits annuellement par la Direction des Affaires Budg\'etaires du Minist\`ere
de la Justice, pour les exercices \textbf{2021 \`a 2026}. Chaque fichier
correspond \`a une situation arr\^et\'ee~: ex\'ecution close pour 2021--2025,
ex\'ecution partielle pour 2026 (donn\'ee provisoire \`a la date du travail).

Chaque fichier comporte plusieurs centaines de lignes correspondant aux
articles budg\'etaires du minist\`ere selon la nomenclature LOF~130-13~:
\textbf{Chapitre} (1211 fonctionnement personnel, 1212 fonctionnement
mat\'eriel, 1213 investissement), \textbf{Programme}, \textbf{R\'egion},
\textbf{Projet} et \textbf{Ligne budg\'etaire} (\textit{Lb}). La cl\'e
compos\'ee \texttt{Chap-Prog-Reg-Proj-Lb}, not\'ee dans la suite
\texttt{Line\_Key}, identifie de mani\`ere unique une rubrique de d\'epense.

Pour chaque ligne et chaque ann\'ee, les colonnes utilis\'ees sont~:

\begin{itemize}
    \item \textbf{CO} : cr\'edits ouverts (montant initial + virements re\c{c}us).
    \item \textbf{Engage\_Reel} : total engag\'e en fin d'exercice (ou \`a date pour 2026).
    \item \textbf{Taux\_Reel} = \texttt{Engage\_Reel}\,/\,\texttt{CO} : taux d'engagement annuel.
\end{itemize}

\subsubsection*{4.1.2 \quad Cha\^ine de pr\'eparation}
\addcontentsline{toc}{subsubsection}{4.1.2 \quad Cha\^ine de pr\'eparation}

L'extraction est r\'ealis\'ee par l'outil HTML autonome d\'ecrit au
chapitre~5. Elle se d\'eroule en cinq \'etapes~:

\begin{enumerate}
    \item \textbf{Lecture des six fichiers \texttt{.xls}} via la biblioth\`eque
    SheetJS, directement dans le navigateur du gestionnaire~;
    \item \textbf{Normalisation des intitul\'es} (suppression des doubles
    espaces, harmonisation des accents et des codes num\'eriques) afin que les
    cl\'es \texttt{Line\_Key} soient strictement comparables d'une ann\'ee \`a
    l'autre~;
    \item \textbf{Calcul des montants nets} (engagements bruts moins
    annulations \'eventuelles)~;
    \item \textbf{Pivot annuel} pour reconstituer, ligne par ligne, la
    s\'erie temporelle des taux d'engagement sur 2021--2026~;
    \item \textbf{Filtrage} des lignes pr\'esentes sur l'ensemble de la p\'eriode
    et dont les cr\'edits ouverts ne sont jamais nuls (un \texttt{CO} nul rend
    le taux d'engagement non d\'efini).
\end{enumerate}

\subsubsection*{4.1.3 \quad P\'erim\`etre final}
\addcontentsline{toc}{subsubsection}{4.1.3 \quad P\'erim\`etre final}

Au terme de la cha\^ine, le jeu de donn\'ees consolid\'e couvre~:

\begin{itemize}
    \item \textbf{107~lignes budg\'etaires communes} \`a tous les exercices 2021--2026~;
    \item r\'eparties sur les trois chapitres principaux du budget du minist\`ere
    (1211 personnel, 1212 mat\'eriel et d\'epenses diverses, 1213
    investissement)~;
    \item soit \textbf{642 observations annuelles} (107 lignes $\times$ 6
    ann\'ees) sur la profondeur historique brute.
\end{itemize}

Apr\`es application du protocole de backtest (section~4.3) qui n\'ecessite
plusieurs ann\'ees ant\'erieures pour entra\^iner les mod\`eles, le jeu
exploit\'e pour la s\'election compte \textbf{321~pr\'evisions \'evalu\'ees}~--
chacune correspondant \`a une ligne donn\'ee test\'ee sur une ann\'ee
donn\'ee (107~lignes $\times$ 3~ann\'ees de backtest 2023--2025).

%================================================================
\phantomsection
\subsection*{4.2 \quad Statistiques descriptives}
\addcontentsline{toc}{subsection}{4.2 Statistiques descriptives}
%================================================================

\subsubsection*{4.2.1 \quad Structure du portefeuille budg\'etaire}
\addcontentsline{toc}{subsubsection}{4.2.1 \quad Structure du portefeuille budg\'etaire}

Sur la p\'eriode 2021--2025 (ann\'ees pleinement cl\^otur\'ees), le
portefeuille consolid\'e des 107~lignes repr\'esente un volume cumul\'e
de plusieurs milliards de dirhams de cr\'edits ouverts, pour un volume
engag\'e l\'eg\`erement inf\'erieur. Le \textbf{taux d'ex\'ecution moyen}
sur la p\'eriode s'\'etablit aux environs de \textbf{70\,\%}, en ligne
avec la moyenne observ\'ee dans les rapports d'ex\'ecution publi\'es par
la Tr\'esorerie G\'en\'erale du Royaume.

Ce taux global masque toutefois une h\'et\'erog\'en\'eit\'e marqu\'ee~: certaines
lignes ex\'ecutent moins de~20\,\% de leur dotation (typiquement, des projets
d'investissement dont les march\'es publics ont \'echou\'e), tandis que d'autres
d\'epassent largement 100\,\% par recours \`a des virements de cr\'edits en
cours d'ann\'ee.

\subsubsection*{4.2.2 \quad Diversit\'e des comportements budg\'etaires}
\addcontentsline{toc}{subsubsection}{4.2.2 \quad Diversit\'e des comportements budg\'etaires}

L'analyse de la \emph{stabilit\'e} (\'ecart-type du taux d'engagement par ligne
sur l'historique 2021--2025) r\'ev\`ele trois r\'egimes structurellement
diff\'erents~:

\begin{itemize}
    \item \textbf{Lignes stables}~: \'ecart-type inf\'erieur \`a 5~points de
    pourcentage. Typiquement, les d\'epenses de personnel (traitements,
    indemnit\'es), o\`u la rigidit\'e statutaire impose une ex\'ecution proche
    du budget~;
    \item \textbf{Lignes mod\'er\'ement volatiles}~: \'ecart-type entre 5 et
    15~points. Souvent~: certaines d\'epenses de fonctionnement courantes
    (fournitures, entretien)~;
    \item \textbf{Lignes volatiles}~: \'ecart-type sup\'erieur \`a 15~points.
    Concentre les d\'epenses d'investissement, les redevances (\'electricit\'e,
    carburants), et les march\'es ponctuels. Pour ces lignes, le taux
    d'engagement peut osciller entre 30\,\% et 200\,\% d'une ann\'ee sur
    l'autre.
\end{itemize}

Cette tripartition justifie a posteriori le choix de la section~3.3 de ne pas
imposer un mod\`ele unique \`a toutes les lignes~: aucun pr\'edicteur ne peut,
de mani\`ere uniforme, traiter aussi bien une masse salariale rigide qu'une
ligne de redevances soumise \`a des pics impr\'evisibles. Le protocole de
s\'election ligne par ligne (section~4.3) est la r\'eponse op\'erationnelle
\`a cette h\'et\'erog\'en\'eit\'e.

%================================================================
\phantomsection
\subsection*{4.3 \quad Protocole exp\'erimental}
\addcontentsline{toc}{subsection}{4.3 Protocole exp\'erimental}
%================================================================

\subsubsection*{4.3.1 \quad Backtests historiques glissants}
\addcontentsline{toc}{subsubsection}{4.3.1 \quad Backtests historiques glissants}

Conform\'ement au cadre m\'ethodologique pos\'e en~3.3.4, la s\'election du
meilleur mod\`ele s'appuie sur un \textbf{backtest \`a horizon glissant} sur
trois ann\'ees de validation, $y \in \{2023, 2024, 2025\}$. Pour chaque ann\'ee
de validation~:

\begin{enumerate}
    \item les cinq mod\`eles (tendance lin\'eaire, moyenne, moyenne pond\'er\'ee,
    persistance, lissage exponentiel simple) sont entra\^in\'es sur les ann\'ees
    strictement ant\'erieures \`a~$y$~;
    \item chaque mod\`ele produit une pr\'ediction $\widehat{\tau}_y$ du taux
    d'engagement~;
    \item l'erreur en valeur absolue relative est calcul\'ee~:
    $\text{MAPE}_y = |\widehat{\tau}_y - \tau_y|\,/\,|\tau_y| \times 100$~;
    \item pour chaque ligne, le mod\`ele retenu est celui dont la moyenne du
    MAPE sur les trois ann\'ees de validation est la plus faible.
\end{enumerate}

Ce protocole pr\'esente deux propri\'et\'es importantes. D'une part, il
n'\'evalue jamais un mod\`ele sur les m\^emes ann\'ees que celles ayant servi
\`a l'entra\^iner~-- contrainte fondamentale pour pr\'evenir le sur-apprentissage.
D'autre part, il accepte de juger chaque mod\`ele sur \emph{trois r\'ealisations
ind\'ependantes}, ce qui att\'enue le bruit d'une mauvaise ann\'ee isol\'ee.

\subsubsection*{4.3.2 \quad G\'en\'eration des projections 2026--2028}
\addcontentsline{toc}{subsubsection}{4.3.2 \quad G\'en\'eration des projections 2026--2028}

Une fois le meilleur mod\`ele identifi\'e pour chaque ligne, il est r\'eentra\^in\'e
sur \emph{l'ensemble} des ann\'ees disponibles (2021--2025, voire 2026 pour les
lignes dont l'ex\'ecution est d\'ej\`a renseign\'ee), puis utilis\'e pour
produire la pr\'evision officielle des ann\'ees 2026, 2027 et 2028. Les
intervalles de confiance \`a $\pm 2\sigma$ sont calcul\'es \`a partir des
r\'esidus du mod\`ele sur l'historique de la ligne.

%================================================================
\phantomsection
\subsection*{4.4 \quad Performance des mod\`eles de pr\'evision}
\addcontentsline{toc}{subsection}{4.4 Performance des mod\`eles de pr\'evision}
%================================================================

\subsubsection*{4.4.1 \quad Comparaison des familles de mod\`eles}
\addcontentsline{toc}{subsubsection}{4.4.1 \quad Comparaison des familles de mod\`eles}

Le tableau~\ref{tab:model-counts} pr\'esente le nombre de pr\'evisions
(une ligne $\times$ une ann\'ee de test) o\`u chaque mod\`ele a \'et\'e s\'electionn\'e comme meilleur sur
les 321~observations du portefeuille triennal.

\begin{table}[h]
\centering
\caption{R\'epartition des mod\`eles s\'electionn\'es sur le portefeuille des 107~lignes (321 pr\'evisions \'evalu\'ees, soit 107~lignes $\times$ 3~ann\'ees de backtest).}
\label{tab:model-counts}
\begin{tabular}{lrr}
\toprule
\textbf{Mod\`ele s\'electionn\'e} & \textbf{Nombre} & \textbf{Part} \\
\midrule
Tendance lin\'eaire (\textit{linear\_trend})        & 132 & 41{,}1\,\% \\
Moyenne historique (\textit{average})               & 93  & 29{,}0\,\% \\
Persistance (\textit{naive})                        & 57  & 17{,}8\,\% \\
Lissage exponentiel simple (\textit{exp\_smoothing})& 30  & 9{,}3\,\%  \\
Moyenne pond\'er\'ee (\textit{weighted\_avg})       & 9   & 2{,}8\,\%  \\
\midrule
\textbf{Total} & \textbf{321} & \textbf{100\,\%} \\
\bottomrule
\end{tabular}
\end{table}

Trois constats m\'eritent d'\^etre soulign\'es. \textbf{Premier constat~:} aucun
mod\`ele ne capte plus de 41\,\% des situations. L'h\'et\'erog\'en\'eit\'e
empirique du portefeuille~-- d\'ej\`a anticip\'ee par la triple stabilit\'e
de~4.2.2~-- est donc confirm\'ee par les r\'esultats du backtest. Tout choix
d'un mod\`ele unique \og pour tout le minist\`ere \fg{} aurait laiss\'e
inexploit\'e plus de la moiti\'e des configurations possibles.

\textbf{Deuxi\`eme constat~:} la tendance lin\'eaire et la moyenne historique,
les deux mod\`eles les plus parcimonieux du jeu (un et z\'ero param\`etre
respectivement, hors estimation de l'intercept), \emph{ensemble} captent
\textbf{70\,\%} des situations. Le compromis biais-variance favorise donc bien,
empiriquement, les mod\`eles les plus simples~-- conform\'ement \`a l'argument
th\'eorique de la section~3.3.2.

\textbf{Troisi\`eme constat~:} le lissage exponentiel simple, dont la
litt\'erature acad\'emique souligne r\'eguli\`erement la robustesse, ne capte
ici que 9\,\% des situations. Son param\`etre de lissage~$\alpha$,
m\^eme cal\'e par optimisation, est p\'enalis\'e par l'\'echantillon
court~: cinq points ne permettent pas de distinguer de mani\`ere fiable une
dynamique exponentielle d'un simple bruit autour d'une moyenne.

\subsubsection*{4.4.2 \quad Analyse des erreurs de pr\'evision}
\addcontentsline{toc}{subsubsection}{4.4.2 \quad Analyse des erreurs de pr\'evision}

La distribution du MAPE de validation sur les 258~pr\'evisions \'evalu\'ees (lignes test\'ees sur 2023--2025)
disposant d'une valeur calcul\'ee est tr\`es asym\'etrique~:

\begin{itemize}
    \item m\'ediane = \textbf{28{,}3\,\%}~;
    \item premier quartile = 0\,\% (lignes parfaitement pr\'edites par
    persistance ou moyenne historique constante)~;
    \item troisi\`eme quartile = 51{,}2\,\%~;
    \item maximum = 2825\,\% (lignes \og pathologiques \fg{} dont le taux
    d'engagement varie de plusieurs ordres de grandeur).
\end{itemize}

Appliqu\'es \`a ces MAPE, les seuils de fiabilit\'e d\'efinis en~3.5.1
($\leq$\,20\,\% \og Fiable \fg{}, 20--50\,\% \og Acceptable \fg{}, 50--100\,\%
\og \`A valider manuellement \fg{}, $>$\,100\,\% \og Incertain \fg{}) produisent
la r\'epartition pr\'esent\'ee au tableau~\ref{tab:reliability}.

\begin{table}[h]
\centering
\caption{R\'epartition du portefeuille par niveau de fiabilit\'e du mod\`ele s\'electionn\'e (321 pr\'evisions \'evalu\'ees).}
\label{tab:reliability}
\begin{tabular}{lrr}
\toprule
\textbf{Niveau de fiabilit\'e} & \textbf{Nombre} & \textbf{Part} \\
\midrule
Fiable                       & 111 & 34{,}6\,\% \\
\`A valider manuellement     & 99  & 30{,}8\,\% \\
Acceptable                   & 75  & 23{,}4\,\% \\
Incertain                    & 36  & 11{,}2\,\% \\
\midrule
\textbf{Total} & \textbf{321} & \textbf{100\,\%} \\
\bottomrule
\end{tabular}
\end{table}

En lecture op\'erationnelle~: \textbf{58\,\%} des situations
(Fiable + Acceptable) peuvent \^etre utilis\'ees directement par le
gestionnaire pour alimenter la phase de cadrage~; \textbf{31\,\%} m\'eritent
une v\'erification au cas par cas avant int\'egration~; \textbf{11\,\%}
restent en zone d'incertitude et doivent \^etre trait\'ees par dire d'expert.
Cette gradation~-- explicitement transmise au gestionnaire dans les
livrables du chapitre~5~-- est l'apport principal du dispositif par rapport
\`a une projection \og brute \fg{} sans qualification.

\paragraph{MAPE pond\'er\'e par les montants.}
La m\'ediane par ligne masque l'impact financier des erreurs. Le
\emph{MAPE pond\'er\'e}, d\'efini comme $\sum_i |\widehat{y}_i - y_i|
\,/\, \sum_i |y_i|$, mesure l'\'ecart en montant agr\'eg\'e plut\^ot
qu'en pourcentage par ligne. Sur le m\^eme backtest 2023--2025, il
s'\'el\`eve \`a \textbf{38{,}1\,\%}, valeur sup\'erieure \`a la m\'ediane.
Cette diff\'erence traduit le fait que les erreurs absolues les plus
\'elev\'ees se concentrent sur quelques grandes lignes (notamment les
r\'emun\'erations) dont l'\'ecart relatif, bien que mod\'er\'e, p\`ese
fortement en valeur. Cette information est restitu\'ee telle quelle au
gestionnaire~: la pr\'ecision varie selon la ligne, et le tableau de
bord n'incite pas \`a un usage m\'ecanique de la pr\'evision globale.

\paragraph{R\'ef\'erence \`a la litt\'erature.}
Cette performance se situe dans la fourchette typique observ\'ee
en pr\'evision sur s\'eries annuelles courtes. Les comp\'etitions de
r\'ef\'erence~\citep{Makridakis2018} rapportent des MAPE m\'edians
de l'ordre de 25 \`a 40\,\% sur des r\'egimes \`a faible profondeur
historique, ce qui place le dispositif dans une plage attendue
plut\^ot que dans une situation d'\'echec.

\paragraph{Distribution bimodale et performance conditionnelle.}
La m\'ediane globale agr\`ege en r\'ealit\'e deux r\'egimes distincts
qu'il convient de dissocier. En appliquant un filtre de volatilit\'e
sur l'historique~2020--2024 (coefficient de variation $\sigma/|\mu|<5\,\%$),
\textbf{30 lignes sur 107} apparaissent structurellement constantes~--
typiquement des subventions reconduites \`a l'identique ou des dotations
fixes~-- pour lesquelles la \og pr\'evision \fg{} se r\'eduit \`a une
reproduction de la valeur ant\'erieure et atteint m\'ecaniquement un
MAPE proche de z\'ero (m\'ediane $0{,}0\,\%$ sur les 59 pr\'evisions
correspondantes du backtest). Sur les \textbf{77 lignes volatiles}
restantes~-- celles o\`u la pr\'evision constitue un v\'eritable
exercice~-- la m\'ediane conditionnelle s'\'el\`eve \`a \textbf{38{,}6\,\%},
contre \textbf{48{,}7\,\%} pour la persistance na\"ive sur le m\^eme
sous-ensemble. C'est sur ce sous-portefeuille, et non sur l'ensemble,
que le dispositif d\'emontre une valeur ajout\'ee~: il r\'eduit l'erreur
m\'ediane d'environ dix points de pourcentage l\`a o\`u la r\'eponse
n'est pas triviale. Cette d\'ecomposition est explicitement assum\'ee~:
la performance globale annonc\'ee ne doit pas \^etre interpr\'et\'ee
comme un succ\`es uniforme.

\subsubsection*{4.4.3 \quad Limites et port\'ee de la pr\'evision budg\'etaire}
\addcontentsline{toc}{subsubsection}{4.4.3 \quad Limites et port\'ee de la pr\'evision budg\'etaire}

Avant de retenir les cinq m\'ethodes int\'egr\'ees au dispositif, plusieurs
approches plus complexes ont \'et\'e \'evalu\'ees lors de la phase
exploratoire (section~3.3.1). Toutefois, dans un contexte caract\'eris\'e
par des historiques tr\`es courts, ces mod\`eles n'ont pas apport\'e
d'am\'elioration significative par rapport aux m\'ethodes plus simples
et ont donc \'et\'e \'ecart\'es.

Au-del\`a de cette consid\'eration m\'ethodologique, la pr\'evision
budg\'etaire pr\'esente une limite inh\'erente \`a son domaine
d'application. Une partie de l'\'evolution des d\'epenses d\'epend de
d\'ecisions futures qui ne sont pas observables dans les donn\'ees
historiques~: arbitrages politiques, nouvelles priorit\'es strat\'egiques,
cr\'edits exceptionnels inscrits en loi de finances ou encore
r\'eallocations d\'ecid\'ees en cours d'exercice. Ces \'ev\'enements
rel\`event de choix de gestion et ne peuvent \^etre anticip\'es de
mani\`ere fiable par un mod\`ele statistique ou algorithmique.

Cette limite ne remet pas en cause l'int\'er\^et du dispositif~; elle
permet au contraire d'en pr\'eciser le r\^ole. L'objectif n'est pas de
remplacer le jugement du gestionnaire, mais de le compl\'eter. Le
syst\`eme automatise l'analyse des lignes dont le comportement historique
pr\'esente une certaine r\'egularit\'e et fournit des estimations
accompagn\'ees d'un niveau de confiance. En parall\`ele, il identifie
les lignes atypiques ou instables n\'ecessitant une analyse approfondie
par les responsables budg\'etaires.

Ainsi, la principale valeur ajout\'ee du dispositif ne r\'eside pas
uniquement dans la production d'une pr\'evision, mais \'egalement dans
sa capacit\'e \`a distinguer les situations pouvant \^etre trait\'ees
de mani\`ere automatis\'ee de celles qui requi\`erent une expertise
humaine. Les indicateurs de fiabilit\'e, de stabilit\'e et les m\'ecanismes
de d\'etection d'anomalies constituent \`a cet \'egard des outils
essentiels d'aide \`a la d\'ecision.

%================================================================
\phantomsection
\subsection*{4.5 \quad Pr\'evision triennale 2026--2028}
\addcontentsline{toc}{subsection}{4.5 Pr\'evision triennale 2026--2028}
%================================================================

\subsubsection*{4.5.1 \quad Vue d'ensemble}
\addcontentsline{toc}{subsubsection}{4.5.1 \quad Vue d'ensemble}

L'application du protocole conduit \`a une programmation triennale agr\'eg\'ee dont le total s'\'etablit \`a quelques milliards de dirhams par exercice, avec une l\'eg\`ere d\'ecroissance apparente sur l'horizon 2026--2028.

Cette d\'ecroissance refl\`ete avant tout la
prudence statistique des mod\`eles parcimonieux~: en l'absence d'un signal
clair de hausse ou de baisse, ils s'alignent sur la moyenne pond\'er\'ee de
l'historique et tendent \`a propager des taux d'engagement inf\'erieurs \`a
100\,\% sur les lignes volatiles. La pr\'evision finale doit donc \^etre lue
comme une \textbf{borne basse statistiquement consistante}, qui sera
n\'ecessairement r\'evis\'ee \`a la hausse par les services lorsque sont
ajout\'es les projets nouveaux et les besoins ponctuels (qui, par construction,
ne figurent pas dans l'historique des cinq ann\'ees pass\'ees).

\subsubsection*{4.5.2 \quad Intervalles de confiance}
\addcontentsline{toc}{subsubsection}{4.5.2 \quad Intervalles de confiance}

Chaque pr\'evision est syst\'ematiquement accompagn\'ee de ses bornes \`a
$\pm 2\sigma$ propres au mod\`ele s\'electionn\'e. Pour les lignes stables,
les intervalles sont \'etroits~: pour la ligne \og Traitements, salaires et
indemnit\'es permanentes du personnel \fg{} (la plus volumineuse du budget),
la largeur relative de l'intervalle de confiance \`a $\pm 2\sigma$ reste
de l'ordre de quelques pourcents autour de la pr\'evision centrale.

Pour les lignes volatiles, les intervalles sont au contraire tr\`es larges et
\emph{informatifs en eux-m\^emes}~: ils signalent au gestionnaire qu'aucune
pr\'ediction ponctuelle ne peut \^etre prise au s\'erieux sans hypoth\`eses
suppl\'ementaires sur la programmation de l'exercice.

%================================================================
\phantomsection
\subsection*{4.6 \quad D\'etection d'anomalies et priorisation des risques}
\addcontentsline{toc}{subsection}{4.6 D\'etection d'anomalies et priorisation des risques}
%================================================================

\subsubsection*{4.6.1 \quad Score Z annuel}
\addcontentsline{toc}{subsubsection}{4.6.1 \quad Score Z annuel}

L'application du score~Z (seuil $|z| > 1{,}5$, voir~3.4.1) sur les 606~observations
(une ligne $\times$ une ann\'ee, sur 2021--2026) fait ressortir \textbf{58 anomalies annuelles}, soit
9{,}6\,\% des observations. Leur r\'epartition annuelle, pr\'esent\'ee au
tableau~\ref{tab:z-per-year}, est relativement uniforme~: aucune ann\'ee ne
concentre plus du quart des signaux, ce qui sugg\`ere que le dispositif capte
bien des \og anomalies de ligne \fg{} et non un effet temporel global (par
exemple, un \'eventuel choc COVID-19 en 2021).

\begin{table}[h]
\centering
\caption{Anomalies annuelles d\'etect\'ees par le score~Z par ann\'ee d'observation.}
\label{tab:z-per-year}
\begin{tabular}{lcccccc}
\toprule
\textbf{Ann\'ee} & 2021 & 2022 & 2023 & 2024 & 2025 & 2026 \\
\midrule
\textbf{Nb anomalies}     & 13   & 6    & 12   & 8    & 7    & 12   \\
\bottomrule
\end{tabular}
\end{table}

\paragraph{Port\'ee et limites du seuil retenu.}
Le score~Z appliqu\'e \`a un historique de cinq \`a six observations par
ligne ne constitue pas un test statistique au sens strict~: l'\'ecart-type
est estim\'e sur un nombre d'observations trop r\'eduit pour qu'un seuil
ponctuel (1{,}5 ou 2) ait une port\'ee probabiliste rigoureuse, et la
m\'ediane attendue d'un d\'etecteur \`a $\pm 1{,}5\sigma$ se situe
m\'ecaniquement autour de 10 \`a 13\,\% des observations sous une
hypoth\`ese de bruit gaussien. Le taux observ\'e de 9{,}6\,\% doit donc
\^etre lu non comme une mesure d'\og anomalit\'e absolue \fg{}, mais comme
le volume de signaux qu'un seuil mod\'er\'ement s\'electif fait remonter
sur ce portefeuille~: il s'agit d'une \emph{heuristique de pr\'e-filtrage},
destin\'ee \`a identifier les \'ecarts notables par rapport \`a la
trajectoire propre de chaque ligne. La pertinence op\'erationnelle du
dispositif ne repose pas sur le score~Z pris isol\'ement, mais sur son
crois\^ement avec l'Isolation Forest (\S\,4.6.2) et sur la matrice de
priorit\'e P0--P4 (\S\,4.6.3), qui filtre les signaux de~Z par leur
coh\'erence avec un profil structurellement atypique. Une ligne signal\'ee
par le seul score~Z reste class\'ee en P1 ou P2, jamais en P0~: le statut
qui lui est attribu\'e est explicitement celui d'une alerte pr\'eventive,
non d'un constat av\'er\'e.

Cette approche ne vise donc pas \`a remplacer l'analyse du gestionnaire
ni \`a r\'ev\'eler des anomalies totalement inconnues. Sa contribution
principale consiste \`a automatiser le pr\'e-filtrage du portefeuille
budg\'etaire, \`a objectiver les crit\`eres d'alerte et \`a hi\'erarchiser
les lignes n\'ecessitant une analyse approfondie. Le dispositif permet
ainsi de concentrer l'attention du gestionnaire sur les situations les
plus pertinentes plut\^ot que sur l'examen manuel de l'ensemble des lignes.

\subsubsection*{4.6.2 \quad Isolation Forest}
\addcontentsline{toc}{subsubsection}{4.6.2 \quad Isolation Forest}

L'Isolation Forest est entra\^in\'e sur les 107~lignes en utilisant comme
vecteur de caract\'eristiques les six taux d'engagement annuels 2021--2026.
Avec un taux de contamination cible de~10\,\% (param\'etrage justifi\'e
en~3.4.2), l'algorithme identifie \textbf{11~lignes au profil structurellement
atypique} (score d'anomalie n\'egatif), parmi lesquelles~:

\begin{itemize}
    \item \og Frais d'h\'ebergement et de restauration \fg{} -- pic
    isol\'e \`a 5097\,\% du cr\'edit ouvert en 2022, contre des taux
    inf\'erieurs \`a 1000\,\% les autres ann\'ees~;
    \item \og Redevances d'\'electricit\'e \fg{} -- taux d'engagement
    syst\'ematiquement sup\'erieur \`a 100\,\% des cr\'edits ouverts (jusqu'\`a
    2500\,\% en 2024), refl\'etant des dotations initiales structurellement
    insuffisantes compens\'ees par virements~;
    \item \og Achat de carburants et lubrifiants \fg{} -- taux moyen
    sup\'erieur \`a 280\,\%~;
    \item \og Maintenance et conservation des archives \fg{} -- taux moyen
    autour de 240\,\%~;
    \item plusieurs lignes de travaux d'am\'enagement, de mat\'eriel
    informatique et de transport ponctuel pr\'esentant des trajectoires
    discontinues sur 2021--2026.
\end{itemize}

Le caract\`ere multivari\'e de l'IF est ici d\'eterminant~: aucune de ces
lignes n'aurait \'et\'e signal\'ee par un simple seuil sur le taux d'une
ann\'ee isol\'ee~; ce sont leurs trajectoires compl\`etes qui les distinguent
du reste du portefeuille.

\paragraph{Stabilit\'e du r\'esultat sous variation du tirage al\'eatoire.}
L'algorithme Isolation Forest reposant sur un \'echantillonnage al\'eatoire
des arbres, on peut l\'egitimement s'interroger sur la sensibilit\'e des
lignes flaggu\'ees au choix de la graine pseudo-al\'eatoire, en particulier
sur un portefeuille de taille mod\'er\'ee ($n=107$). Pour le v\'erifier,
l'algorithme a \'et\'e r\'e-entra\^in\'e \textbf{100~fois} avec des graines
diff\'erentes (param\'etrage par ailleurs identique~: 200~arbres,
contamination~10\,\%). La distribution des fr\'equences de signalement par
ligne est tr\`es asym\'etrique~: \textbf{8~lignes sont flaggu\'ees dans
100\,\% des tirages}, \textbf{9 dans au moins 90\,\%}, \textbf{11~dans au
moins 70\,\%}~--~soit exactement le noyau publi\'e ci-dessus~--~tandis que
\textbf{89~lignes ne sont jamais flaggu\'ees} et que seules 2~lignes se
situent dans une zone \og borderline \fg{} (fr\'equences de 13 et 34\,\%).
Le c\oe ur du diagnostic est donc essentiellement d\'eterministe sur ce
portefeuille, et la d\'ependance r\'esiduelle \`a la graine n'affecte qu'un
nombre marginal de lignes au caract\`ere atypique d\'ej\`a peu marqu\'e.
Cette robustesse, conjointe \`a la r\`egle de croisement avec le score~Z
discut\'ee en~\S\,4.6.3, garantit que la matrice de priorit\'e P0--P4 ne
d\'epend pas d'un al\'ea informatique.

\subsubsection*{4.6.3 \quad Articulation Z $\times$ IF et matrice de priorit\'e}
\addcontentsline{toc}{subsubsection}{4.6.3 \quad Articulation Z $\times$ IF et matrice de priorit\'e}

Le croisement des deux signaux (Z r\'ecent~: anomalie d\'etect\'ee sur les
deux derni\`eres ann\'ees~; IF~: profil structurellement atypique) produit la
matrice de contingence du tableau~\ref{tab:z-if-cross}, sur les
103~lignes\footnote{Le score Risque\_Combine n'est calcul\'e que pour les
lignes dont les six ann\'ees 2021--2026 sont effectivement renseign\'ees~;
4~lignes au profil partiellement renseign\'e en sont exclues, ce qui ram\`ene
le p\'erim\`etre de la matrice P0--P4 \`a 103~lignes sur les 107 du
portefeuille.} pour lesquelles les deux d\'etecteurs produisent un signal
exploitable.

\begin{table}[h]
\centering
\caption{Croisement des deux d\'etecteurs sur les 103~lignes pour lesquelles les deux signaux sont disponibles.}
\label{tab:z-if-cross}
\begin{tabular}{lcc}
\toprule
 & \textbf{IF atypique = Non} & \textbf{IF atypique = Oui} \\
\midrule
\textbf{Z r\'ecent = Non} & 70 & 6 \\
\textbf{Z r\'ecent = Oui} & 22 & 5 \\
\bottomrule
\end{tabular}
\end{table}

L'application de la matrice de priorit\'e P0--P4 d\'efinie en~3.4.3 conduit
\`a la r\'epartition pr\'esent\'ee au tableau~\ref{tab:p0-p4}.

\begin{table}[h]
\centering
\caption{R\'epartition des 103~lignes selon la priorit\'e d'attention recommand\'ee.}
\label{tab:p0-p4}
\begin{tabular}{llrr}
\toprule
\textbf{Code} & \textbf{Label} & \textbf{Nombre} & \textbf{Part} \\
\midrule
P0 & Critique & 5   & 4{,}9\,\%  \\
P1 & \'Elev\'e & 22 & 21{,}4\,\% \\
P2 & Mod\'er\'e & 6 & 5{,}8\,\%  \\
P4 & Normal & 70   & 68{,}0\,\% \\
\midrule
\textbf{Total} & & \textbf{103} & \textbf{100\,\%} \\
\bottomrule
\end{tabular}
\end{table}

Les \textbf{5~lignes P0} (Critique) cumulent les deux signaux~: anomalie Z
r\'ecente \emph{et} profil IF atypique. Elles re\c{c}oivent la recommandation
\og Revue manag\'eriale imm\'ediate \fg{}~: il s'agit notamment des
redevances d'\'electricit\'e, de la ligne \og Achat de carburants et
lubrifiants \fg{} et de quelques lignes de travaux ponctuels.

Les \textbf{22~lignes P1} (\'Elev\'e) pr\'esentent une alerte Z r\'ecente
sans profil structurellement atypique~: ce sont typiquement des lignes qui
ont connu un d\'epassement ponctuel en 2025 ou 2026 mais qui restent dans
le r\'egime g\'en\'eral du portefeuille (par exemple~: \og Frais de
maintenance du mat\'eriel informatique \fg{} avec un taux 2026 sup\'erieur
\`a 110\,\%, ou \og Indemnit\'es de d\'eplacement \fg{} en d\'epassement
mod\'er\'e). La recommandation associ\'ee est \og Renforcer le suivi
mensuel \fg{}.

Les \textbf{6~lignes P2} (Mod\'er\'e) pr\'esentent un profil IF atypique mais
sans alerte Z r\'ecente~: elles d\'erogent au portefeuille moyen, mais leur
trajectoire ne d\'evie pas anormalement par rapport \`a leur propre
historique. La recommandation est \og Surveillance trimestrielle \fg{}.

Les \textbf{70~lignes P4} (Normal, environ 68\,\% du p\'erim\`etre scor\'e)
ne d\'eclenchent aucune alerte et b\'en\'eficient du suivi p\'eriodique
standard. C'est sur ces lignes que les m\'ecanismes habituels de la
circulaire annuelle du Premier ministre suffisent.

\textbf{Lecture op\'erationnelle.} Sur les 103~lignes scor\'ees,
\textbf{33~lignes} (P0 + P1 + P2) requi\`erent une attention soutenue --
chiffre report\'e au tableau de bord sous l'intitul\'e \og Lignes \`a
risque \fg{}. Cette concentration -- environ un tiers du portefeuille --
d\'efinit le \textbf{rayon d'action prioritaire} du contr\^oleur de
gestion~: c'est une fraction du portefeuille suffisamment ramass\'ee pour
qu'un suivi mensuel manuel soit envisageable, mais suffisamment riche pour
couvrir l'essentiel du risque d'ex\'ecution budg\'etaire.

\paragraph{Interpr\'etation op\'erationnelle de la matrice de priorit\'e.}
La matrice P0--P4 doit \^etre interpr\'et\'ee comme un m\'ecanisme de
priorisation des comportements atypiques et non comme un classement des
lignes selon leur poids budg\'etaire. Les niveaux de priorit\'e r\'esultent
du croisement entre des signaux statistiques (score~Z) et structurels
(Isolation Forest), dont l'objectif est d'identifier les lignes dont le
comportement s'\'ecarte de leur trajectoire habituelle ou de celui du
reste du portefeuille.

Par cons\'equent, une ligne repr\'esentant un faible montant budg\'etaire
peut \^etre class\'ee prioritaire si elle pr\'esente une anomalie marqu\'ee,
tandis qu'une ligne de poids financier important peut ne g\'en\'erer aucune
alerte lorsque son ex\'ecution demeure r\'eguli\`ere et pr\'evisible. Cette
logique est coh\'erente avec l'objectif du dispositif, qui consiste \`a
attirer l'attention du gestionnaire sur les situations inhabituelles
n\'ecessitant une analyse compl\'ementaire.

L'\'evaluation de l'impact financier potentiel d'une anomalie rel\`eve
ensuite du jugement du gestionnaire et peut \^etre appr\'eci\'ee \`a
l'aide des informations budg\'etaires disponibles dans le tableau de
bord Power~BI.

\paragraph{Valeur ajout\'ee du dispositif de d\'etection.}
Il convient de souligner que les anomalies identifi\'ees par le dispositif
ne sont pas n\'ecessairement inconnues des gestionnaires budg\'etaires.
Dans la pratique, le suivi r\'egulier des taux d'engagement et des cr\'edits
consomm\'es permet d\'ej\`a de rep\'erer certaines situations atypiques.
La contribution du syst\`eme r\'eside toutefois moins dans la d\'ecouverte
d'anomalies in\'edites que dans l'\textbf{automatisation, l'objectivation
et la priorisation} de leur analyse. D'une part, l'ensemble du portefeuille
budg\'etaire est examin\'e de mani\`ere syst\'ematique selon des crit\`eres
homog\`enes, sans d\'ependre exclusivement de l'exp\'erience individuelle
du gestionnaire. D'autre part, les anomalies d\'etect\'ees sont
hi\'erarchis\'ees selon leur niveau de priorit\'e, ce qui facilite
l'identification rapide des lignes n\'ecessitant une attention particuli\`ere.

L'int\'er\^et op\'erationnel du dispositif est donc de transformer un
processus d'analyse essentiellement manuel en un m\'ecanisme d'aide \`a
la d\'ecision capable de signaler automatiquement les situations \`a
examiner. Dans cette perspective, la valeur ajout\'ee ne r\'eside pas
uniquement dans la d\'etection des anomalies, mais \'egalement dans la
r\'eduction du temps d'analyse et dans la standardisation du processus
de suivi budg\'etaire.

%================================================================
\phantomsection
\subsection*{4.7 \quad Synth\`ese et limites}
\addcontentsline{toc}{subsection}{4.7 Synth\`ese et limites}
%================================================================

\subsubsection*{4.7.1 \quad Synth\`ese}
\addcontentsline{toc}{subsubsection}{4.7.1 \quad Synth\`ese}

Les r\'esultats empiriques de ce chapitre confirment, sur le cas concret des
107~lignes du Minist\`ere de la Justice, la validit\'e du dispositif
m\'ethodologique pos\'e au chapitre~3~:

\begin{itemize}
    \item l'h\'et\'erog\'en\'eit\'e des lignes interdit tout choix d'un mod\`ele
    unique~; la s\'election ligne par ligne est statistiquement justifi\'ee~;
    \item les mod\`eles les plus parcimonieux (tendance lin\'eaire + moyenne)
    captent 70\,\% des situations, validant a posteriori le rejet des
    approches riches en param\`etres~;
    \item la d\'etection d'anomalies focalise l'attention du gestionnaire sur
    33~lignes prioritaires sur 107, dont 5~lignes critiques~;
    \item chaque pr\'evision est qualifi\'ee par un niveau de Confiance Globale
    transmis avec le chiffre, conform\'ement \`a l'exigence de transparence pos\'ee
    en~3.5.
\end{itemize}

\subsubsection*{4.7.2 \quad Limites}
\addcontentsline{toc}{subsubsection}{4.7.2 \quad Limites}

\paragraph{Temporalit\'e d'utilisation du dispositif:}
Les r\'esultats pr\'esent\'es dans ce chapitre reposent sur des
exercices cl\^otur\'es afin de permettre la validation du dispositif.
Cette d\'emarche est conforme aux pratiques usuelles en pr\'evision, qui
n\'ecessitent de comparer les sorties du mod\`ele \`a des observations
connues avant toute utilisation op\'erationnelle.

En situation r\'eelle, le syst\`eme est destin\'e \`a \^etre ex\'ecut\'e
au d\'ebut de chaque nouvel exercice budg\'etaire. Les donn\'ees
historiques disponibles sont alors utilis\'ees pour produire des
pr\'evisions et identifier les lignes n\'ecessitant une attention
particuli\`ere pour l'ann\'ee \`a venir. Les alertes g\'en\'er\'ees n'ont
donc pas vocation \`a r\'einterpr\'eter un exercice d\'ej\`a cl\^otur\'e,
mais \`a orienter le suivi budg\'etaire futur.

Le dispositif doit ainsi \^etre compris comme un outil de pr\'eparation
et de priorisation destin\'e \`a accompagner le gestionnaire dans le
pilotage des exercices \`a venir, et non comme un m\'ecanisme d'analyse
r\'etrospective uniquement.

\paragraph{Cadence annuelle et alignement m\'etier:}
Le choix d'une ex\'ecution annuelle du dispositif ne r\'esulte pas d'une
contrainte technique, mais de la nature des donn\'ees disponibles et du
processus de pr\'evision budg\'etaire observ\'e au sein de la Direction
du Budget. Les exercices de programmation et de pr\'eparation
budg\'etaire s'appuient principalement sur des historiques annuels
consolid\'es, structur\'es selon le calendrier de la loi de finances.

Dans ce contexte, le dispositif a \'et\'e con\c{c}u pour produire des
pr\'evisions et des alertes destin\'ees \`a accompagner les travaux de
planification et de pr\'eparation des exercices futurs. Son horizon
d'analyse est donc coh\'erent avec la temporalit\'e des d\'ecisions
auxquelles il est destin\'e \`a contribuer. Le syst\`eme est ainsi
align\'e sur les pratiques de pr\'evision et de programmation
budg\'etaire observ\'ees au sein de la Direction du Budget.

\medskip

Cinq limites doivent \^etre signal\'ees \`a l'utilisateur~:

\textbf{Premi\`erement}, le caract\`ere descriptif des mod\`eles. Les cinq
mod\`eles retenus prolongent les r\'egularit\'es pass\'ees~; ils
n'int\`egrent ni les d\'ecisions de programmation discr\'etionnaires des
services (nouveau projet de construction d'un palais de justice, par
exemple), ni les chocs exog\`enes (r\'eforme l\'egislative, restructuration
de la carte judiciaire). Pour ces \'ev\'enements, la pr\'evision statistique
doit \^etre compl\'et\'ee par un ajustement manuel document\'e.

\textbf{Deuxi\`emement}, la profondeur historique reste structurellement
courte. Avec cinq exercices clos, les statistiques de fiabilit\'e et de
stabilit\'e ont une marge d'erreur intrins\`eque non n\'egligeable. Le
dispositif est con\c{c}u pour \^etre r\'eentra\^in\'e chaque ann\'ee, et
sa qualit\'e s'am\'eliorera m\'ecaniquement avec l'allongement de la
profondeur.

\textbf{Troisi\`emement}, la d\'etection d'anomalies signale des
\emph{r\'egularit\'es atypiques}, non des fraudes ou des erreurs comptables.
Une ligne signal\'ee P0 m\'erite une revue manag\'eriale, mais celle-ci peut
parfaitement conclure que l'atypicit\'e refl\`ete une r\'ealit\'e
op\'erationnelle l\'egitime (par exemple, des redevances syst\'ematiquement
sous-budg\'etis\'ees par convention administrative). Le dispositif est un
outil d'aide \`a la d\'ecision, jamais un substitut au jugement professionnel.

\textbf{Quatri\`emement}, l'\textit{absence de variables exog\`enes}.
Le dispositif de pr\'evision propos\'e repose exclusivement sur l'historique
propre \`a chaque ligne budg\'etaire. Il n'int\`egre pas explicitement de
variables exog\`enes susceptibles d'influencer l'\'evolution future des
d\'epenses, telles que l'inflation, les prix de l'\'energie, les d\'ecisions
salariales nationales ou les indices sectoriels. Ce choix r\'esulte
principalement de la nature du portefeuille \'etudi\'e~: les 107~lignes
budg\'etaires analys\'ees couvrent des cat\'egories de d\'epenses tr\`es
h\'et\'erog\`enes, dont les d\'eterminants externes diff\`erent fortement.
\`A titre d'exemple, les d\'epenses de carburant peuvent \^etre influenc\'ees
par les cours internationaux du p\'etrole, les d\'epenses de personnel par
les d\'ecisions relatives aux r\'emun\'erations publiques, tandis que
certaines d\'epenses d'investissement peuvent d\'ependre de l'\'evolution
des co\^uts du secteur du b\^atiment et des travaux publics.

L'int\'egration de variables explicatives adapt\'ees \`a chacune des
lignes aurait n\'ecessit\'e la construction et la maintenance d'un grand
nombre de mod\`eles sp\'ecifiques, r\'eduisant significativement la
simplicit\'e, la robustesse et l'op\'erationnalit\'e du dispositif. Le
choix retenu a donc \'et\'e de privil\'egier une approche g\'en\'erique,
homog\`ene et directement applicable \`a l'ensemble du portefeuille
budg\'etaire. Cette limite implique que le syst\`eme ne peut pas anticiper
les ruptures li\'ees \`a des d\'ecisions futures ou \`a des chocs exog\`enes
qui ne figurent pas encore dans les donn\'ees historiques. Les pr\'evisions
produites doivent ainsi \^etre interpr\'et\'ees comme une projection
fond\'ee sur les comportements observ\'es, destin\'ee \`a assister le
gestionnaire dans son analyse, et non comme une anticipation exhaustive
de l'ensemble des facteurs susceptibles d'affecter l'ex\'ecution
budg\'etaire.

\textbf{Cinqui\`emement}, la \textit{validit\'e de la continuit\'e des
lignes budg\'etaires}. L'analyse repose sur l'hypoth\`ese que les lignes
budg\'etaires retenues conservent une signification suffisamment stable
sur l'ensemble de la p\'eriode \'etudi\'ee. Afin de limiter les risques
de rupture historique, seules les lignes pr\'esentant une correspondance
coh\'erente selon la cl\'e composite
Chapitre--Programme--R\'egion--Projet--Libell\'e ont \'et\'e conserv\'ees
dans le p\'erim\`etre final. Cette approche r\'eduit le risque de
confondre des lignes distinctes au cours du temps. Toutefois, comme dans
tout syst\`eme budg\'etaire, certaines \'evolutions organisationnelles
ou modifications de nomenclature peuvent affecter la comparabilit\'e
parfaite des observations historiques. Cette limite doit \^etre prise
en compte dans l'interpr\'etation des r\'esultats et constitue une piste
d'am\'elioration pour de futurs travaux, notamment par l'int\'egration
de tables de correspondance historiques valid\'ees par les gestionnaires
m\'etiers.

\subsubsection*{4.7.3 \quad Transition vers le chapitre 5}
\addcontentsline{toc}{subsubsection}{4.7.3 \quad Transition vers le chapitre 5}

Les r\'esultats num\'eriques pr\'esent\'es dans ce chapitre n'ont d'int\'er\^et
op\'erationnel que s'ils sont accessibles, lisibles et r\'eutilisables par les
gestionnaires du minist\`ere. Le chapitre~5 d\'ecrit l'architecture du
double livrable qui mat\'erialise cet acc\`es~: un outil HTML autonome
ex\'ecut\'e dans le navigateur du gestionnaire, qui produit les rapports
\`a partir des fichiers \texttt{.xls} bruts~; et un tableau de bord Power~BI
qui consomme ces rapports pour offrir une vue de pilotage \`a la Direction
du Budget.
